\chapter{Введение}
\label{chap:intro}
\setlength{\parskip}{0.8em}
\section{Актуальность и мотивация исследования}
Ожоговые травмы представляют собой наиболее серьезную глобальную проблему общественного здравоохранения. По данным Всемирной организации здравоохранения, примерно 11 миллионов человек в мире ежегодно получает ожоги, из которых около 180~000 случая завершаются летально~\cite{who2018burns}. Особенно широкое распространение ожоговых травм наблюдается в странах со средним и низким уровнем дохода, где ожоги составляют существенную долю всех обращений в экстренные службы этих стран.

Все ожоги делятся по глубине поражения тканей на три типа (или степени). Ожоги первой степени, которые также принято называть поверхностными, затрагивают исключительно эпидермис и проявляются покраснением, болью и незначительным отеком. Ожоги второй степени, именуемые также частичной толщины, проникают в дерму, вызывая образование пузырей и выраженную боль. Поскольку ожоги третьей степени (полной толщины) разрушают все слои кожи, включая нервные окончания, больному обязательно требуется хирургическое вмешательство. Точная и своевременно проведённая классификация степени ожога имеет критически важное значение при выборе оптимальной тактики лечения - в частности, ожоги первой степени могут подлежать амбулаторному лечению, ожоги второй степени могут потребовать госпитализации, в то время как ожоги третьей степени требуют интенсивного лечения, включающего в себя впоследствии пересадку кожи. я трансплантацию живой кожи. Однако даже опытнейшие врачи иногда не могут правильно определить визуально степень ожога, тем более может и не хватить стремления, особенно в первые часы после травмы~\cite{monstrey2008assessment}. Искусственный интеллект, представляя собой глубокое обучения и компьютерного зрения, дарует качественно новый путь диагностике: автоматический анализ изображений позволяет оценивать степень ожога без учета субъективных факторов и ускоряет решение. Казахстан, страна с высокой степенью промышленной развитости и быстро растущей системой здравоохранения, остро нуждается в доступных простым людям и эффективных методах диагностики ожоговой травмы. Внедрение в клиническую практику систем, построенному на искусственном интеллекте, могло бы значительно улучшить качество медицинской помощи, особенно в отдалённых её уголках. \section{Постановка задачи}

Клинической задачей, рассматриваемой в диссертации, может быть названа автоматическая морфологическая классификация степени ожогового поражения по фотоснимкам. Примечанием к вышеописанному следует добавить, что данное исследование вызвано необходимостью преодоления трех соображений, представляющих собой одноразовые ограничения, а именно

\begin{enumerate}
\item Данная сложность диагностики --- визуальная оценка ожога зависит от опыта врача, и от времени прошедшего с момента травмы, и от условий освещения~\cite{monstrey2008assessment}. \item Доступность специализированной экспертизы --- далеко не все медицинские учреждения, особенно расположенные в сельской местности, обладают возможностью обращаться к специалистам-комбустиологам. \item Темп диагностирования --- в неотложной ситуации требуется максимально оперативная предварительная оценка, обеспечивающая сортировку по степени тяжести и возможность планирования тактики лечения~\cite{peck2012burns}. \end{list}

В диссертации была предложена система по детекции объектов, основанная на архитектуре YOLOv8, позволяющая решать все три упомянутые проблемы: проводить объективную оценку на стандартной фотографии, не требовать специального оборудования и функционировать в реальном времени. \section{Цели и задачи исследования}

Наиболее приоритетной задачей является задача разработки, последующего обучения и валидирования модели глубокого обучения, которая будет способны автоматически классифицировать степени ожогов. Определить перед собой одно из основных испытаний:

\begin{enumerate}
\item Осуществить основательный систематизирующий обзор актуальных методов компьютерного зрения по анализу кожных поражений и ожогов.
\end{enumerate} \item Подготовить и проанализировать датасет Kaggle Skin Burn Dataset, оценив распределение классов и оценив качество аннотаций.
\item Реализовать конвейер предобработки данных с аугментацией для повышения обобщающей способности обученной модели. \item Обучить на подготовленном датасете модель YOLOv8n, применив трансферное обучение.
\item Оценить производительность модели по метрикам mAP, precision и recall.
\item Разработать REST API для практического применения модели в клинических условиях. \end{enumerate}

\section{Научная новизна}

Научная новизна работы состоит из четырех аспектов.

1. YOLOv8 и его применение в задаче классификации ожогов. Современная система детекции объектов в реальном времени, основанная на архитектуре YOLOv8~\cite{ultralytics2023yolov8}. Она применяется для классификации ожогов. В отличие от тех старых добрых стандартных CNN-классификаторов~\cite{he2016resnet,simonyan2014vgg}, YOLO не только выполняет задачу классификации изображения, но еще и при этом, параллельно выделяет зону поражения на нем.

2.- Построение конвейера обработки медицинских изображений. Спроектирован и реализован унифицированный конвейер предобработки (включающий такие способы аугментации данных~\cite{shorten2019augmentation} как HSV-преобразования, масштабирование, зеркальное отображение) нормализацию и формирование аннотаций, подходящих для YOLO.

3. Метрики детекции объектов для диагностики в медицине. Для оценки качества качества обнаруженных поражений, используем метрику mAP (mean Average Precision) с порогами IoU, унаследованную от стандартных бенчмарков~\cite{lin2014coco,everingham2010pascal} и адаптированную для наших задач:

\begin{equation}
\text{mAP} = \frac{1}{|\mathcal{C}|} \sum_{c \in \mathcal{C}} \text{ \begin{equation}
\text{AP}_c
\label{eq:map}
\end{equation}

\noindent где $\mathcal{C}$ - соотношение классов (степеней ожогов); $\text{AP}_c$ - средняя точность для класса~$c$, подсчитываемая как площадь под кривой precision-recall.

4. Практическая реализация системы с REST API. А вкорне, по сравнению с чисто исследовательскими работами, настоящая диссертация содержит генерацию полностью рабочей и готовой к внедрению разработки REST API на базе фреймворка Litestar~\cite{litestar2023}, а разрешение интеграции модели в имеющиеся медицинские информационные системы. \section{Практическая значимость}

Практическими результатами, вносящими вклад в реализацию данной работы:

\begin{itemize}
\item Обученная модель YOLOv8n для классификации трех степеней ожогов с весами, справедливыми для повторного использования. \item Реализуемый конвейер данных на Питоне (Ultralytics, PyTorch) с документированной системой подготовки датасета, обучения и инференса. \item REST API, который предназначен для синхронизации с клиническими информационными системами, который поддерживает, в числе прочего, обработку изображений в распространенном формате JPEG и возвращает детализированные и структурированные результаты детекции.
\item Спредлагаемая методология обучения модели на платформах облачного графического процессинга (vast). ai), относительно применимая для прочих задач медицинской визуализации.
\item Готовая концепция для интеграции в медицинские учреждения Казахстана в качестве вспомогательного средства диагностики.
\end{itemize}

\section{Структура диссертации}

Диссертация состоит из пяти глав. Глава~\ref{chap:literature} содержит обзор литературы по глубокому обучению в медицинской визуализации и эволюции архитектур YOLO. Глава~\ref{chap:methodology} описывает набор данных, конвейер предобработки, архитектуру YOLOv8 и фреймворк оценки. Глава~\ref{chap:results} представляет результаты экспериментов, метрики производительности модели и примеры детекции. Глава~\ref{chap:conclusion} синтезирует вклады, обсуждает клинические импликации и предлагает направления будущих исследований.
